\documentclass[11pt]{article}
\usepackage{amsmath,amssymb,amsthm,mathtools,hyperref}

\begin{document}
\section*{Normal tissues carrying oncogenic mutations}
\label{prob:250}
\label{rank:38}
\par\smallskip

\noindent\textbf{Domain.} Clone-to-neoplasia risk boundaries \quad \textbf{Retained status/source context.} Phenotypically normal tissues with selected cancer-associated mutations are established observations. The unresolved programme is explaining and predicting progression in tissue context; a universal sharp boundary is not established.
\paragraph{Definitions and assumptions.} For a specified adult tissue let $G$ be observed alteration categories, $C$ normalized clone size and $E$ measured tissue/host context. Let $H_0$ indicate independently normal baseline histology, $T_N$ first qualifying neoplasia and $T_D$ death before neoplasia. A finite-dimensional competing-risk law specifies cause-specific hazards $\lambda_N(t\mid G,C,E)$ and $\lambda_D(t\mid G,C,E)$ plus any shared frailty distribution. These hazards refer to the baseline-normal population $H_0=1$ and the risk set still alive and free of qualifying neoplasia at time $t$. Their displayed notation suppresses this eligibility and risk-set conditioning. A concrete instance must state whether the hazards are marginal over latent frailty or conditional on a declared frailty variable $Z$. In the latter representation, specify both hazards conditional on $Z$ and the distribution of $Z$ conditional on $H_0=1,G,C,E$; obtain the target probability by averaging the conditional cumulative incidence over that distribution. No independence of hypothetical latent event times is assumed merely by calling the law a competing-risk law. Observation error and clone-to-lesion linkage uncertainty are declared. Normal-tissue expansion and malignant progression are distinct. Biopsies sample only part of the tissue; missed small clones are explicitly modeled.
\paragraph{Problem statement.} Find a calibrated competing-risk progression probability and supported risk boundary:
\[q_H(G,C,E)=\mathbb{P}(T_N\leq H,T_N<T_D\mid H_0=1,G,C,E),\qquad
B_H=\{(G,C,E):q_H(G,C,E)\geq q_0\}.\]
Determine whether $B_H$ is identifiable from existing normal/tumour and longitudinal observations, with uncertainty and sensitivity to unmeasured tissue context and competing mortality. A threshold $q_0\in(0,1)$ is a chosen risk criterion, not evidence of a biological separatrix. The two marginal cause-specific hazards determine this cumulative-incidence function; when frailty-conditional hazards are used instead, their specification and the baseline conditional frailty distribution jointly determine it.
\paragraph{Additional binding assumptions and resolution criterion.} Fix the prediction time origin, baseline normal-histology criterion and clone-sampling process. If clone size and tissue state evolve, either condition on baseline covariates in a declared reduced hazard model or supply their joint time-dependent law. A risk superlevel set is an operational classification region, not a deterministic transition surface. Report competing-risk uncertainty and lack of support for unseen clone/context combinations.
\paragraph{Outstanding instance specifications.} Before fitting or testing, declare the adult target population, tissue compartment and prediction unit. If clones are indexed, specify their selection rule and how multiple clones in one host share outcomes. State whether a qualifying lesion must descend from the indexed clone or may be any qualifying lesion in a declared tissue or host. Define the qualifying histological endpoint and whether $T_N$ denotes biological onset, first detectability or first diagnosis; specify absent-event and event-tie conventions consistently with the strict inequality $T_N<T_D$. These endpoint and attribution choices are not selected by this scheme. Specify alteration categories and co-occurrence rules, the clone-size denominator and scale, measured context variables and units, and the relation between measured proxies and latent quantities. Declare the horizon $H$ and its units, the chosen $q_0$, and the admissible finite-dimensional model family under the binding catalogue conventions. Identification is relative to declared observation designs and data resources, including applicable entry, selection, surveillance, censoring, interval observation, death ascertainment and follow-up management mechanisms, and assumptions relating their observed laws to the target law. Predeclare the estimator, calibration and informative predictive-performance criteria, uncertainty assessment and validation domain; calibration alone does not establish informative prediction. State whether identification of $B_H$ is required modulo covariate-null sets in the baseline-normal population or pointwise on a declared domain with assumptions supporting that claim. This equality convention remains an instance choice. Distinguish identification of the threshold region from identification of the entire probability function; uncertainty in $q_H$ need not change every threshold decision.

\paragraph{Alternative formulations.} (i) Normal-versus-tumour cross-sectional comparison identifies distributional differences but not prospective progression hazards. (ii) Longitudinal clone follow-up estimates risk under adequate linkage and ascertainment. (iii) A mechanistic stochastic first-passage model asks a stronger dynamical question; stable mutant-normal states cannot be inferred solely from the absence of diagnosed cancer in a short record.
\paragraph{Obstructions and partial progress.} Martincorena et al. show selected mutant clones in aging normal human oesophagus, demonstrating that cancer-associated alterations need not immediately imply malignancy. The exact challenge seeks the tissue-level constraints and transition process. Clonal fitness, mutation identity and malignant potential therefore cannot be collapsed into one score. Rare progression, spatial sampling and uncertain clone ancestry limit identification; an informative result may be a bound rather than one universal transition surface.
\paragraph{Sources (supplied provenance; not all rechecked).}
\begin{itemize}
\item Cancer Grand Challenges, active challenge index, accessed 11 October 2026. Research-agenda evidence, not a mathematical conjecture. \url{https://www.cancergrandchallenges.org/challenges/active-challenges}
\item Cancer Grand Challenges, Normal phenotypes; active tissue-level programme. \url{https://www.cancergrandchallenges.org/challenges/active-challenges/normal-phenotypes}
\item I. Martincorena et al., Somatic mutant clones colonize the human esophagus with age (2018); established normal-tissue mutation/selection observations. \url{https://pmc.ncbi.nlm.nih.gov/articles/PMC6298579/}
\end{itemize}


\section*{Binding catalogue conventions}
This edition revises both supplied catalogues and preserves original problem numbers 1--330. It adds one globally unique label \verb|\label{rank:x}| per problem, where smaller $x$ means greater provisional difficulty. A rank label is an anchor at that problem's section: \verb|\ref{rank:x}| returns the section number, not the literal difficulty number. The visible \texttt{rank:x} text and the index give the difficulty number itself.

The entries are \emph{parameterized mathematical research schemes}, including empirical prediction, causal identification, historical reconstruction, theory and engineering. They are not 330 certified unsolved mathematical conjectures. A rigorous instance fixes the objects and restrictions below; changing those objects changes the question. Established identities, finite calculations, non-identification constructions and solved special cases are not being posed anew as open. Some residual tasks concern informative bounds, biological realization or external validation rather than a new theorem.

Every main definition and problem statement received a mathematical editorial review. Each entry now has an additional binding-conditions paragraph; direct corrections are separately recorded in the accompanying audit. The alternative formulations inherit the applicable conventions, but are different targets unless equivalence is proved. Where an old broad phrase conflicts with an explicit restriction below it, the explicit restriction controls.

\textbf{Literature boundary.} Thirteen selective primary-literature checks cover eighteen entries in this revision; their access levels and scope are recorded explicitly. The supplied bibliography and its earlier source-access descriptions are retained as provenance, not represented as newly verified. Neither this review nor a source's description proves the absence of later solutions. The date of this edition is 13 October 2026; there is no claim of exhaustive current-status certification for all 330 entries.

\subsection*{Binding conventions for a concrete instance}
\paragraph{Objects, quantifiers and non-vacuity.}
An instance declares measurable state, input, output and parameter spaces; a fixed candidate family; admissible initial laws, interventions and environments; a finite horizon or specified asymptotic regime; observation and sampling laws; and performance metrics, positive normalization scales and positive tolerances. These are chosen before fitting or testing, not adjusted after seeing failures. Shared parameters must genuinely be shared across the declared validation domain. Finite dimension alone does not restrict arbitrary real-number encodings: specify a finite library, bounded description length, or regularity/complexity bounds. Exclude constant or no-signal solutions through a fixed informative-output or baseline-improvement requirement where appropriate. A universally quantified result cannot be established by choosing one trivial favorable instance.

\paragraph{Probability laws and data.}
Every expectation or probability refers to the specified stochastic law, including initial uncertainty, process noise, observation noise and sampling as appropriate. A deterministic equation is not a probability model without a declared random input or initial law. ``Observed'' targets in population-level formulas mean population observables $g_j^*$, not noisy estimates asserted to be exact. Let $\widetilde g_j$ be independent validation estimates and suppose
\[
\mathbb{P}\{\lvert\widetilde g_j-g_j^*\rvert\le r_j\text{ for every }j\}\ge1-\alpha.
\]
Then an uncertainty-aware sufficient certification rule for positive target tolerances $\varepsilon_j$ is
\[
\max_j\frac{|\widehat g_j-\widetilde g_j|+r_j}{\varepsilon_j}\le1.
\]
On that simultaneous confidence event it implies every population error is at most its tolerance, by the triangle inequality. Vector/norm versions use corresponding valid simultaneous bounds. Training uncertainty must also be included when the claim averages over repeated training data. Evaluation on finitely many conditions does not establish a continuous-domain supremum without covering and regularity bounds.

\paragraph{Dynamics and well-posedness.}
For ODEs and SDEs give coefficient domains, initial laws and existence/uniqueness conditions in the intended solution sense, or keep the entire admitted solution set. For stochastic jump models specify nonnegative off-diagonal rates, diagonal normalization, rate orientation and non-explosion. For PDEs specify spatial domains, boundary/initial conditions, conservation or source terms, and a weak or classical solution class. Moving-domain concentration equations include dilution/transport terms. A stationary law is used only for a time-homogeneous augmented model under existence assumptions; a stationary approximation to a driven system must be justified. Absorbing extinction states are retained. Survival-conditioned, quasi-stationary and unconditional laws are distinct and must not be substituted for one another.

\paragraph{Identification and causal meaning.}
Two parameter values are observationally equivalent when they induce the same full distribution under all allowed observation/intervention designs, not merely the same likelihood value for one sample. A target is identified only if constant on the equivalence class. Gauge symmetries, label switching and reference conventions are quotiented or fixed. Finite-data compatible sets use a declared confidence, likelihood or discrepancy rule; an empty set signals model rejection. Bayesian credible sets are not automatically confidence sets. A prior or regularizer does not create data identification. Bounds intended to resolve a question must be informative at a predeclared width or decision threshold.

A causal contrast uses well-defined intervention versions and a stated causal model. Random assignment, or explicit consistency, overlap and exchangeability assumptions, is required for causal interpretation from observations. Interference and time-varying confounding are represented when present. Fixing a mediator by intervention differs from conditioning on its observed value. Marginal potential-outcome laws do not determine individual cross-world effects or an arbitrary joint coupling. Ethically admissible designs, rather than unconstrained interventions, define the experimental menu.

\paragraph{Metrics, units, optimization and limits.}
Specify the sampling or reference measure for densities and integrals, the log base, and dimensionless arguments of logarithms. Total variation is $\sup_A|P(A)-Q(A)|$; Wasserstein distances require the declared ground metric and appropriate moment conditions; KL divergence uses its stated absolute-continuity convention. All normed readouts are scaled consistently, and error probabilities use $0<\alpha<1$. A ratio has a positive denominator or an explicit undefined/zero convention. Thermodynamic entropy production is nonnegative in expectation under the relevant hypotheses; individual stochastic realizations need not be. An infimum is not an attained minimum without compactness/coercivity and continuity or another existence argument. Fix the order of large-population, long-time and vanishing-mutation limits. Stability claims declare perturbation space and symmetries.

\subsection*{How the provisional difficulty order was assigned}
Four editorial components are scored from 1 (lower burden) to 5 (higher burden): (a) formal or inferential difficulty, (b) access to evidence and identification burden, (c) mechanistic/computational integration, and (d) independent validation and transfer. Their unweighted total orders the entries. For selected entries in the highest bands, an explicit editorial priority breaks equal totals; all other ties are broken by original problem number. Equal totals are \emph{not} evidence of reliably different true difficulties. Scores are judgments, not measurements, completion-time estimates or consensus rankings.

The object ranked is an informative, nondegenerate resolution of the \emph{scoped} scientific target, with the stated budget and validation domain. A sharp useful impossibility result can qualify; an already known identity, vacuous bound, arbitrary favorable model or lookup table cannot. Rank 330 is merely the least difficult relative assessment in this catalogue, not a claim of an easy biological solution. Broad historical titles do not automatically outrank precisely formulated mechanistic tasks.



\end{document}
